% SPDX-License-Identifier: MIT
\chapter{Komponenten, Widgets und Formulare}\label{chap:widgets}
\ChapterSummary{Dieses Kapitel fuehrt durch die gebrauchsfertigen
Grund-Controls, ihre Zustandseigenschaften und ihre Lebensdauer.
Sie erfahren, wann ein Widget genuegt und wann eine nichtvisuelle
Komponente oder ein Modell die bessere Wahl ist.}

\section{Ein Besitzbaum, zwei Sichtweisen}
\term{Component} ist die Grundlage fuer Objekte, die eine
einheitliche Eigentums- und Lebensdauerbeziehung benoetigen.
\term{Widget} erweitert dieses Konzept um Sichtbarkeit,
Aktivierbarkeit, Fokusfaehigkeit, Groesse, Position und Ereignisse.
\term{Container} besitzt Komponenten ueber eindeutiges
\texttt{std::unique\_ptr}-Eigentum. Nur Widgets bilden dabei
sichtbare Kinder; eine Command-Komponente im gleichen
Container erscheint nicht im Layout.

\begin{lstlisting}[language=SASDCpp,caption={Eigentum ausdruecklich uebergeben}]
Window window;
auto& form = window.emplace<VBox>();
auto& caption = form.emplace<Label>("Einstellungen");
auto& save = form.emplace<Button>("Speichern");
\end{lstlisting}
Der \texttt{Window}-Destruktor entsorgt die verschachtelten
Komponenten. Fuer dynamische Strukturen existieren
\texttt{adopt(std::unique\_ptr<Component>)} und
\texttt{release(Component\&)}, womit die Lebensdauer bewusst
uebertragen werden kann. Ein nach \texttt{release} entfernter
Teilbaum muss im Anwendungscode entweder weiterbesessen oder
zerstoert werden.

\Important{C++-Referenzen auf Widgets sind nicht selbst
lebensdauersicher. Wenn Sie einen Container aufloesen oder
ein Kind entfernen, duerfen spaetere Callbacks nicht
mehr ueber alte Referenzen zugreifen.}

\section{Gemeinsame Widget-Eigenschaften}
Die wichtigsten Methoden aus \filepath{include/sasd/ui/widget.hpp}:
\begin{tabularx}{\linewidth}{@{}p{39mm}X@{}}
\toprule
Methode & Bedeutung\\\midrule
\texttt{setVisible(bool)} & Beeinflusst Darstellung und Teilnahme am Layout.\\
\texttt{setEnabled(bool)} & Sperrt oder erlaubt normale Benutzeraktionen.\\
\texttt{setFocusable(bool)} & Definiert grundsaetzliche Fokusfaehigkeit.\\
\texttt{hasFocus()} & Zeigt aktuellen logischen Fokus an.\\
\texttt{setSizeConstraints(...)} & Begrenzt die gewuenschte Groesse.\\
\texttt{measure(...)} & Berechnet die benoetigte Groesse.\\
\texttt{arrange(Rect)} & Setzt endgueltige Grenzen.\\
\bottomrule
\end{tabularx}
\texttt{canReceiveFocus()} beruecksichtigt
\emph{focusable, visible und enabled} gleichzeitig.
Ein deaktivierter Button sollte daher weder als erreichbares
Tab-Ziel noch als ausfuehrbare Aktion behandelt werden.

\section{Label und Button}
Ein \term{Label} transportiert UTF-8-Text, aber keine
eingebaute Editierfunktion. Ein \term{Button} bietet
\texttt{setText}, \texttt{setOnActivated} und
\texttt{activate}. Ein aktivierbarer Button kann per
Tastatur oder Zeigegeraet ausgeloest werden. Das eigene
Pressed-Feedback wird getrennt vom dauerhaften Enabled-Zustand
gefuehrt.

\begin{lstlisting}[language=SASDCpp]
auto& status = form.emplace<Label>("Noch nichts gespeichert");
auto& save = form.emplace<Button>("Speichern");
save.setOnActivated([&status] {
    status.setText("Speichern angefordert");
});
\end{lstlisting}
Soll derselbe Vorgang auch aus Menue und Shortcut erreichbar sein,
ist ein \term{Command} meist sinnvoller; Kapitel
\ref{chap:commands} zeigt die Bindung.

\section{TextField: Unicode und Auswahl}
\term{TextField} ist ein einzelzeiliger UTF-8-Editor.
Eingaben werden aus \term{TextInputEvent} verarbeitet,
Tastaturbewegung und Loeschen arbeiten mit
Unicode-\emph{Scalars} statt rohen Bytes.
Das ist ein wichtiger Unterschied: Ein mehrbyteiges UTF-8-Zeichen
darf beim Rueckwaertsloeschen nicht zerstoert werden.

Die API kennt \texttt{text()}, \texttt{setText()},
\texttt{cursorPosition()}, \texttt{setCursorPosition()},
\texttt{setSelection()}, \texttt{selectAll()} und
\texttt{selectedText()}. Die Auswahl hat getrennte
Anker- und Cursorpositionen. Es existieren ausserdem
\texttt{copySelectionToClipboard()},
\texttt{cutSelectionToClipboard()} und
\texttt{pasteFromClipboard()} mit einem expliziten
\term{Clipboard}-Objekt.

\begin{lstlisting}[language=SASDCpp]
auto& entry = form.emplace<TextField>("Beispiel");
entry.selectAll();
const std::string selected = entry.selectedText();
entry.setCursorPosition(0);
\end{lstlisting}
Die Cursorposition wird in Unicode-Scalars gemessen,
nicht in Bytes und nicht in Terminalzellen. Komplexe
Graphemcluster koennen mehrere Scalars umfassen.
Breite Zeichen koennen im Terminal mehr als eine Zelle
belegen. Diese Groessen muessen getrennt bleiben.

\section{CheckBox, RadioButton und RadioGroup}
\term{CheckBox} bildet einen booleschen Zustand ab:
\begin{lstlisting}[language=SASDCpp]
auto& logging = form.emplace<CheckBox>("Protokoll", true);
logging.setOnCheckedChanged([](bool checked) {
    // Anwendungseinstellungen aktualisieren.
});
logging.setChecked(false);
\end{lstlisting}
Ein \term{RadioGroup} ist eine \emph{nichtvisuelle}
Komponente fuer gegenseitig ausschliessende Optionen.
RadioButtons referenzieren diese Gruppe explizit.
Sie wird nicht aus benachbarten Widgets geraten.

\begin{lstlisting}[language=SASDCpp]
RadioGroup modeGroup; // Muss die Buttons ueberleben.
Window window;
auto& form = window.emplace<VBox>();
auto& read = form.emplace<RadioButton>(modeGroup,"Nur lesen");
auto& edit = form.emplace<RadioButton>(modeGroup,"Bearbeiten");
read.setSelected(true);
\end{lstlisting}
Die Deklarationsreihenfolge sorgt hier dafuer,
dass beim Verlassen des Scopes zuerst das Fenster
und seine Buttons, danach die Gruppe zerstoert werden.

\section{ComboBox: Auswahl ist nicht Vorschau}
\term{ComboBox} besitzt eine lokale Liste von Strings
und einen optionalen gewaehlten Index. Sie ist derzeit
\emph{nicht editierbar}. Methoden wie
\texttt{setSelectedIndex}, \texttt{selectedText},
\texttt{appendItem} und \texttt{setItems} steuern
die festgeschriebene Auswahl.
Bei Austausch der gesamten Liste wird eine alte
Auswahl bewusst geloescht, damit derselbe Zahlenindex
nicht unbemerkt eine andere Bedeutung bekommt.

Ein geoeffnetes Dropdown hat einen zweiten,
vorlaeufigen \term{Preview-Index}. Navigation veraendert
zunaechst nur die Vorschau. \textbf{Enter} uebernimmt,
\textbf{Escape} verwirft, \textbf{F4} steuert die
Dropdown-Absicht. Das ist fuer Formulare wichtig,
deren teure oder gefaehrliche Folgeaktionen erst
nach einer bestaetigten Wahl laufen sollen.
\begin{lstlisting}[language=SASDCpp]
auto& target = form.emplace<ComboBox>(
    std::vector<std::string>{"Test", "Produktion"});
target.setSelectedIndex(0);
target.setOnSelectionChanged([](auto index) {
    // Nur bestaetigte Auswahl verarbeiten.
});
\end{lstlisting}

\section{Textstil statt plattformspezifischer Farbe}
\term{TextStyle} ermoeglicht eine reduzierte, portable
Textgestaltung mit Farbangaben und Attributen.
\begin{lstlisting}[language=SASDCpp]
TextStyle style;
style.foreground = Color::bright_cyan;
style.bold = true;
caption.setTextStyle(style);
\end{lstlisting}
Die Wirkung ist von den Faehigkeiten des Backends
abhaengig. Verwenden Sie Farbe als Unterstuetzung,
nicht als einziges Traegermedium fuer Bedeutung.
\PracticalTip{Trennen Sie Statusdaten, sichtbaren Text und
Benutzerinteraktion. Ein TextField ist kein Datenmodell;
ein Label ist keine Protokollhistorie; ein Command ist
keine kuenstlich verborgene Schaltflaeche.}
